\section{The weak-activity height recursion}
\label{wh:section}

This section constructs the height recursion on tori and proves its
consequences when a microscopic smallness estimate is available.  Its local
maps also apply to a small activity supplied at a later scale.  The two
assertions have different hypotheses: the microscopic construction below
applies to a fixed interaction at sufficiently large temperature, or to a
sufficiently spread-out interaction at a prescribed coefficient greater than
$8\pi$.  For a general fixed interaction near its physical transition, the
small initial activity will instead be supplied by Theorem~\ref{en:entry}.

Write
\begin{equation}\label{wh:conventions}
 Q=-\Delta_J,\qquad A_0=Q/v^2,\qquad v=v_J,\qquad
 k=\beta/v^2,\qquad k_0=8\pi.
\end{equation}
Write $a(k)=a_J(k)$ for the structural coefficient of
Theorem~\ref{st:free-rectangle-limit}.
We denote the held-fixed reference coefficient by $a'$ and put
$\alpha=\sqrt{a'}$.  Thus $a'$ is a proposed output coefficient, whereas
$k$ is the physical input parameter.  On nonconstant modes use covariance
$A_0^{-1}$, and integrate the constant mode uniformly over
$[0,2\pi/\alpha)$.  Insert the normalized comb
$\sum_{q\in\Z}e^{iq\alpha\phi_x}$ at every vertex and the interaction
\begin{equation}\label{wh:compensation}
 \exp\{s(\phi,A_0\phi)/2\},\qquad s=1-a'/k.
\end{equation}
The resulting measure, with $\sigma=\alpha\phi$, is exactly the height
measure modulo constants: its quadratic exponent is
$-(1-s)(\phi,A_0\phi)/2=-(\sigma,A_0\sigma)/(2k)$.
Comb identities can first be read with positive Fourier approximate identities and then
passed to the limit after the independent Gaussian integration below.
All convolutions are independent of $k$ when $a'$ is held fixed.

For a block $B$ of side $m$ set
\begin{equation}\label{wh:singletons}
 e_B^0(\phi)=\frac1{2|J|v^2}\sum_{x\in B}\sum_{y\in J}
       (\phi_{x+y}-\phi_x)^2,
 \qquad c_B^\alpha(\phi)=m^{-2}\sum_{x\in B}\cos(\alpha\phi_x).
\end{equation}
In particular $\sum_Be_B^0=(\phi,A_0\phi)$, including the displayed
factor $1/2$ accounting for the two orientations of an edge.

The weak-coupling analysis is related to the sine-Gordon contraction
of Dimock--Hurd \citep{DimockHurd2000} and the discrete Gaussian maps
of Bauerschmidt--Park--Rodriguez \citep{BPR2024I}.
The present map, its complex activity analyticity, and its near-marginal
second-order coefficients are proved below with the norms used here.
Their later application to a fixed rough lattice model also requires
the physical entry theorem of Section~\ref{en:section}.

\subsection{Finite-range kernels and the ultraviolet preparation}

We first give the estimates uniformly for
$J_\rho=\{x\in\Z^2\setminus\{0\}:|x|_\infty\le\rho\}$.
Let $d_\rho$ be the largest power of two not exceeding the integer $\rho$,
and set $\gamma=3/5$.  Write $N_\rho=(2\rho+1)^2$ and
\[
 b_\rho(t)=\frac{\sin((2\rho+1)t/2)}{(2\rho+1)\sin(t/2)}.
\]
The exact symbols are
\begin{equation}\label{wh:spread-symbol}
 Q(\xi)=\frac{N_\rho}{N_\rho-1}
          \{1-b_\rho(\xi_1)b_\rho(\xi_2)\},\qquad
 v^2=\frac{N_\rho\rho(\rho+1)}{6(N_\rho-1)}.
\end{equation}
On $[-\pi,\pi]$ one has $|b_\rho(t)|\le1$ and
$|b_\rho(t)|\le\pi/((2\rho+1)|t|)$.  Where $b_\rho$ is negative,
$|t|\ge2\pi/(2\rho+1)$, so $b_\rho\ge-1/2$.
Consequently $0\le\gamma Q\le1-c$ for all sufficiently large $\rho$.
Taylor expansion in a fixed $\rho\xi$ window, and the preceding bound
outside a sufficiently large such window, give
\begin{equation}\label{wh:ellipticity}
 v^2\asymp\rho^2,\qquad
 Q(\xi)\asymp\min\{\rho^2|\xi|^2,1\},\qquad
 Q(\xi)=v^2|\xi|^2\{1+O(\rho^2|\xi|^2)\}.
\end{equation}
For the lower bound in the remaining compact annulus, any subsequential
limit is the product of two sinc functions; its value is strictly less
than one away from the origin.  This proves a uniform positive lower
bound there and completes the ellipticity argument.

Fix a sufficiently large power of two $L$, and let $m_0$ be the first
power of $L$ at least $\rho^B$, where the absolute exponent $B$ is chosen
after the finite list of difference estimates below.  Define
\[
 P_t(\cos z)=\left(\frac{\sin(tz/2)}{t\sin(z/2)}\right)^{64},
 \qquad \mathsf P_t=P_t(1-Q/16).
\]
First take independent normals of variance $\gamma v^2$; then take
\begin{equation}\label{wh:covariances}
 C_{\rm uv}=\frac{v^2(1-\gamma Q)}{Q}(1-\mathsf P_{m_0/d_\rho}),
 \qquad
 \Gamma_m=\frac{v^2(1-\gamma Q)}{Q}
       (\mathsf P_{m/d_\rho}-\mathsf P_{Lm/d_\rho}),
 \quad m=m_0L^j.
\end{equation}
On the torus of side $n=L^N>D_0m_0$, let $m_*$ be the largest member
of $\{m_0L^j:j\ge0\}$ strictly below $n/D_0$.  This set contains
$m_0$.  Use shells whose upper side is at most $m_*$; when
$m_*=m_0$ there are no long shells.  The remaining
nonconstant covariance is
$v^2(1-\gamma Q)\mathsf P_{m_*/d_\rho}/Q$; retain the independent
uniform mean.  Here $D_0$ is fixed after all support radii have been fixed.

\begin{lemma}[Kernel bounds]\label{wh:kernels}
The two kernels in Equation~\eqref{wh:covariances} are positive
semidefinite polynomial kernels, of ranges at most $Cm_0$ and
$CLm$, respectively.  For the long-wave shells $\Gamma_m$ with
$m\ge m_0$, and their binary pieces at the corresponding lower
scales, the pointwise and positive-order operator bounds of
Lemma~\ref{rg:kernels} hold uniformly in sufficiently large $\rho$,
using ordinary nearest-axis differences.  The positive-order
diagonal and operator bounds also hold for the final smoothed
remainder at its stopping scale.  Consequently
Lemma~\ref{rg:regulator} applies to these shell and terminal
integrations.

The ultraviolet covariance is controlled separately by
\begin{equation}\label{wh:uv-energy}
 0\le A_0^{1/2}C_{\rm uv}A_0^{1/2}
   =(1-\gamma Q)(1-\mathsf P_{m_0/d_\rho})\le\Id.
\end{equation}
No scale-$m_0$ nearest-axis derivative or whole-regulator
averaging bound is asserted for $C_{\rm uv}$.

For each fixed admissible scale ratio $R$, the rescaled plane
shell from lower side $m$ to upper side $Rm$, and its scaled
differences through order eight, converge to the corresponding
derivatives of
\begin{equation}\label{wh:continuum-symbol}
 \widehat G_R(\eta)=\frac{P_u(\eta)-P_u(R\eta)}{|\eta|^2},
 \qquad
 P_u(\eta)=\left(\frac{\sin(u|\eta|)}{u|\eta|}\right)^{64},
 \qquad u=\frac{v}{d_\rho\sqrt{32}}.
\end{equation}
The parameter $u$ ranges over a compact subinterval of $(0,\infty)$.
The errors have bounds $C(\rho^C/m)^c$, summable over the starting
scales and tending to zero as $\rho\to\infty$.  In particular, on the
plane, for bounded $\alpha$,
\begin{equation}\label{wh:eigenvalue}
 \lambda_j:=L^2e^{-\alpha^2\Gamma_{m_0L^j}(0,0)/2}
  =\lambda(1+\epsilon_j),\qquad
 \lambda=L^{2-a'/(4\pi)},\qquad
 \sum_{j\ge0}|\epsilon_j|=o(1).
\end{equation}
The ultraviolet diagonal obeys
\begin{equation}\label{wh:uv-diagonal}
 D_{\rm uv}:=\gamma v^2+C_{\rm uv}(0,0)=v^2(1+o(1)).
\end{equation}
For an arbitrary fixed admissible $J$ the shell and terminal long-wave assertions hold with
constants depending on $J$, a sufficiently small $\gamma>0$, and a
sufficiently large starting scale.  In that version the summable errors
tend to zero as the starting scale tends to infinity; neither
Equation~\eqref{wh:uv-diagonal} nor $\gamma>1/2$ is asserted.

For the physical entry construction there is a second, fixed-$J$
family.  With integral $M$ and the same polynomials $\mathsf P_M$, set
\[
 C^{\rm en}_{<M}=\frac{v^2}{Q}(1-\mathsf P_M),\qquad
 \Gamma^{\rm en}_M=\frac{v^2}{Q}(\mathsf P_M-\mathsf P_{LM}).
\]
These are positive semidefinite polynomial kernels of ranges at most
$C_JM$ and $C_JLM$, respectively, and
$0\le A_0^{1/2}C^{\rm en}_{<M}A_0^{1/2}=1-\mathsf P_M\le\Id$.
For $M$ sufficiently large, the plane shells and the torus shells
whose upper side $LM$ is below $n/D_0$, with $D_0$ sufficiently large
depending on $J$, satisfy
\[
 \begin{split}
 |\Gamma^{\rm en}_M(x,y)|
  &\le C_J\left(1+\log^+\frac{LM}{M+\dist(x,y)}\right),
 \qquad \Gamma^{\rm en}_M(x,x)\ge c_J\log L,\\
 |\nabla_x^a\nabla_y^c\Gamma^{\rm en}_M(x,y)|
  &\le C_{J,r}(M+\dist(x,y))^{-r},
       \qquad r=|a|+|c|,\quad 1\le r\le12,\\
 \|\Gamma^{\rm en}_M\|_{2\to2}&\le C_J(LM)^2,\qquad
 \|\nabla^r(\Gamma^{\rm en}_M)^{1/2}\|_{2\to2}^2
       \le C_{J,r}M^{-2(r-1)},\qquad 1\le r\le6 .
 \end{split}
\]
Here the constants are chosen before $L$.  In particular the
positive-order diagonal bounds of Lemma~\ref{rg:kernels} hold.
For the mean-free terminal covariance
$\Gamma^{\rm en}_*=v^2\mathsf P_{M_*}/Q$ on a torus of side
$n\ge2M_*$, those diagonal and positive-order operator bounds hold
at scale $M_*$, with constants independent of $n/M_*$ and $L$.
Its undifferentiated diagonal satisfies
$0\le\Gamma_*^{\rm en}(x,x)\le C_J(1+\log(n/M_*))$.

At fixed $J,L$, the continuum shell for this family has
$\upsilon_{\rm en}=v/\sqrt{32}$ in place of $u$ in
Equation~\eqref{wh:continuum-symbol}.  Writing it as
$G_{L,\upsilon_{\rm en}}$, for $r=|a|+|c|\le8$ one has
\[
 \sup_{x,y\in\Z^2}\left|
 M^r\nabla_x^a\nabla_y^c\Gamma^{\rm en}_M(x,y)
 -\left.
   \partial_X^a\partial_Y^cG_{L,\upsilon_{\rm en}}(X-Y)
  \right|_{X=x/M,\,Y=y/M}\right|
 \le C_{J,L}M^{-1/2}.
\]
For a compact interval $I\Subset(0,\infty)$ of reference coefficients,
the actual fundamental multipliers for a sufficiently large start $m$
are therefore
\[
 \begin{gathered}
 \delta_j^{\rm en}=\Gamma^{\rm en}_{mL^j}(0,0)-\frac{\log L}{2\pi},
 \qquad
 \lambda_j^{\rm en}=L^2e^{-a'\Gamma^{\rm en}_{mL^j}(0,0)/2}
       =\lambda(1+\epsilon_j^{\rm en}),\\
 \epsilon_j^{\rm en}=e^{-a'\delta_j^{\rm en}/2}-1,\qquad
 \sum_{j\ge0}|\epsilon_j^{\rm en}|
 \le \frac{C_{I,J,L}m^{-1/2}}{1-L^{-1/2}}
       \quad(a'\in I).
 \end{gathered}
\]
The shell diagonals and these multipliers are identical on a retained
torus and the plane.  This entry-family assertion is only for fixed
$J$; no uniformity as the interaction range tends to infinity is
asserted.
\end{lemma}

\begin{proof}
For integer scale ratios the elementary identity
\[
 P_{Lt}(\cos z)=P_t(\cos z)
       \left(\frac{\sin(Ltz/2)}{L\sin(tz/2)}\right)^{64}
       \le P_t(\cos z)
\]
holds, with values at removable singularities defined by continuity.
The degree of $P_t$ is $32(t-1)$, and
$P_t(1-q/16)=1-(t^2-1)q/3+O(q^2)$.
Thus the quotients by $Q$ are polynomials with nonnegative spectral
values, also at zero.  A polynomial of degree $d$ in $Q$ has range at
most $d\rho$ in sup distance, proving the range bounds.
Moreover, $0\le1-\gamma Q\le\Id$ and $0\le\mathsf P_t\le\Id$
on the spectrum.  Multiplication by $A_0^{1/2}$ on both sides
of the ultraviolet multiplier proves Equation~\eqref{wh:uv-energy};
both sides of its equality vanish on the constant mode.

For a long-wave shell $\Gamma_m$, split it into its binary pieces.
In $|\xi|\le c/\rho$, the angle
$z=\arccos(1-Q(\xi)/16)$ satisfies
$(m/d_\rho)z\asymp m|\xi|$, and the binary covariance multiplier is
bounded by
\[
 C\min\{m^2,|\xi|^{-2}(m|\xi|)^{-64}\}.
\]
Outside this region it is bounded by
$C\rho^2(\rho/m)^{64}$.  Multiplying by the required nearest-difference
symbols and integrating gives, for a binary piece at side $m$,
$|\nabla^r\Gamma|\le C_rm^{-r}$ for every required $r\le12$.
The high-frequency integral has an extra factor $\rho^{66}m^{-64}$;
increasing $B$ absorbs this factor after each of the finitely many
rescalings.  Spectral supremum estimates give
$\|\Gamma_m\|\le C(Lm)^2$ and
$\|\nabla^r\Gamma_m^{1/2}\|^2\le C_rm^{-2(r-1)}$ for $1\le r\le6$.
For the terminal remainder use the same low- and high-frequency
multiplier bounds, with $\mathsf P_{m_*/d_\rho}$ in place of the
shell difference.  Positive-order spectral suprema and the normalized
torus Fourier sum split at $|\xi|=m_*^{-1}$ give the operator and
diagonal bounds at scale $m_*$ whenever $n\ge2m_*$.  In the low
region the radial annulus count gives $C_rm_*^{-2r}$ for the
diagonal of $r$ differences, $1\le r\le6$; in the high region the
same factor $\rho^{66}m_*^{-64}$ is absorbed by the choice of $B$.
These constants do not require an upper bound for $n/m_*$.
Summing the binary pointwise estimates and using their ranges gives
the logarithmic kernel bound and the positive-order mixed-distance
bounds.  A fixed low-frequency annulus in each
binary piece gives its strictly positive diagonal lower bound.  The
Fourier sums have enough frequencies in those annuli by the choice of
$D_0$.  These are precisely the kernel hypotheses used in
Lemmas~\ref{rg:regulator} and~\ref{rg:charged}.

For convergence of these shells, put $\eta=m\xi$ and divide the integral at
$|\eta|=(m/\rho)^\varepsilon$, with a fixed sufficiently small
$\varepsilon>0$.  The tail just estimated is a negative power of
$m/\rho$, even after eight scaled differences.  In the inner region,
\[
 \frac{m}{2d_\rho}\arccos(1-Q(\eta/m)/16)
 =u|\eta|\{1+O((\rho|\eta|/m)^2)\}.
\]
Taylor expansion of the sine ratios and of the difference multipliers
therefore has a power-small integrated error.  All quotients are regular
at zero, so no additional infrared truncation is required.  On the
diagonal the continuum integral is
\[
 G_L(0)=\frac1{2\pi}\int_0^\infty
       (P_u(r)-P_u(Lr))\frac{\dd r}{r}
       =\frac{\log L}{2\pi}.
\]
The last identity follows by integrating from $\delta$ to $T$, changing
variables in the second term, and sending $\delta\downarrow0$ and
$T\uparrow\infty$, using $P_u(0)=1$ and its decay.  This proves
Equation~\eqref{wh:eigenvalue} with geometric summability.

The multiplier for $D_{\rm uv}/v^2$ tends to one almost everywhere as
$\rho\to\infty$.  Outside $|\xi|<c/\rho$ it is uniformly bounded.
Inside that disk, the bound
$C\min\{m_0^2/\rho^2,1/(\rho^2|\xi|^2)\}$ has integral at most
$C\rho^{-2}(1+\log m_0)$.  Dominated convergence outside, followed by
this vanishing inside contribution, proves
Equation~\eqref{wh:uv-diagonal}.  Range makes the ultraviolet entries
identical on all sufficiently large tori.

For fixed $J$, let $d$ be the largest power of two not exceeding its
sup range, and use $m/d$ in the polynomial indices.  Square symmetry gives
$Q(\xi)=v^2|\xi|^2+O_J(|\xi|^4)$; inclusion of the nearest neighbors
implies that zero is its only Fourier zero.  Choose
$\gamma\sup Q<1$ and repeat the shell and terminal arguments with
the fixed range in place of $\rho$.  The polynomial positivity,
range, and ultraviolet energy bound follow directly as before.
This proves the fixed-$J$ assertions for the first family.

We give the estimates for the entry family with its literal,
unrescaled polynomial indices.  On the Fourier square,
$Q(\xi)\asymp_J|\xi|^2$.  For
$z=\arccos(1-Q(\xi)/16)$ one has exactly
$\sin(z/2)=\sqrt{Q(\xi)/32}$.  The elementary sine bound and the
binary identity give
\[
 \begin{split}
 0\le\mathsf P_t(\xi)&\le C\min\{1,(t\sqrt{Q(\xi)})^{-64}\},\\
 D_t(\xi):=\mathsf P_t-\mathsf P_{2t}
 &=\mathsf P_t\{1-\cos^{64}(tz/2)\}\\
 &\le C\min\{t^2Q(\xi),\,\min(1,(t\sqrt{Q(\xi)})^{-64})\}.
 \end{split}
\]
Consequently the binary covariance
$\Delta_t=v^2D_t/Q$ is bounded by
$C_J\min\{t^2,|\xi|^{-2}\min(1,(t|\xi|)^{-64})\}$.
It is a polynomial of degree $O(t)$, so its range is $C_Jt$.
Multiplication by the nearest-axis difference symbols and radial
integration, split at $|\xi|=t^{-1}$, prove
$|\nabla_x^a\nabla_y^c\Delta_t(x,y)|\le C_{J,r}t^{-r}$
for $0\le r=|a|+|c|\le12$.  Spectral suprema similarly give
$\|\nabla^r\Delta_t^{1/2}\|^2\le C_{J,r}t^{2-2r}$ for $0\le r\le6$.
Summing over the binary sides $t=M,2M,\ldots,LM/2$ proves the
operator bound at order zero and at orders $2\le r\le6$; at order one
use directly
$|\xi|^2v^2(\mathsf P_M-\mathsf P_{LM})/Q\le C_J$.
At positive pointwise order, a piece can contribute only when
$t\ge c_J\dist(x,y)$, up to the fixed difference stencil.
The geometric sum of $t^{-r}$ then gives
$C_{J,r}(M+\dist(x,y))^{-r}$, and the same counting at order
zero gives the displayed logarithm.  On a fixed annulus
$c_J/t\le|\xi|\le C_J/t$ with its constants sufficiently small,
$\mathsf P_t$ is bounded below and
$1-\cos^{64}(tz/2)\ge c_Jt^2Q$ for all large $t$.
Its integral gives a positive constant to each binary diagonal.
This proves the lower bound.  Choosing $D_0$ larger than twice the
polynomial range constant, enlarged for the fixed difference stencils,
makes each retained torus kernel the
periodization of its plane kernel with at most one contributing
image, and makes its diagonal identical to the plane diagonal.
Thus all these bounds hold on the retained tori as stated.

For the terminal covariance, its nonconstant multiplier is bounded by
\[
 C_J|\xi|^{-2}\min\{1,(M_*|\xi|)^{-64}\}.
\]
Its spectral supremum after $2r$ differences proves the positive-order operator
bound.  Splitting the normalized torus Fourier sum at
$|\xi|=M_*^{-1}$ proves the positive-order diagonal bound:
the number of modes in every radial dyadic annulus is bounded by
its area times $Cn^2$, with the same bound for the innermost disk
when $n\ge2M_*$.  For $1\le r\le6$ the resulting sum is at most
$C_{J,r}M_*^{-2r}$.  This argument has no dependence on an upper
bound for $n/M_*$.  At order zero the annuli between the smallest
nonzero frequency and $M_*^{-1}$ each contribute at most $C_J$,
and there are $O(1+\log(n/M_*))$ of them; the higher-frequency
sum is bounded.  This proves the stated terminal value-diagonal
bound.  Polynomial division at zero proves positivity,
range, and the asserted energy inequality for $C^{\rm en}_{<M}$.

For the stated rate of continuum convergence put $\eta=M\xi$ and
use
\[
 Q(\eta/M)=v^2|\eta|^2/M^2+O_J(|\eta|^4/M^4),\qquad
 \frac M2z(\eta/M)=\upsilon_{\rm en}|\eta|
                \{1+O_J(|\eta|^2/M^2)\}.
\]
The rescaled multiplier is
$v^2(\mathsf P_M-\mathsf P_{LM})/(M^2Q)$.
On $|\eta|\le M^{1/100}$, write each sine ratio as
$[z/(2\sin(z/2))]^{64}[\sin(tz/2)/(tz/2)]^{64}$.
Taylor's formula, the preceding two expansions, and the bounds
$|(\sin r/r)^{64}|\le C(1+r)^{-64}$ and
$|((\sin r/r)^{64})'|\le C(1+r)^{-64}$ for $r\ge1$
give an error bounded by
$C_{J,L}M^{-2}(1+|\eta|)^{-60}$ after division by $|\eta|^2$.
At zero the difference of the two sine ratios vanishes quadratically,
so this bound extends across the removable quotient.  The product of
$r$ scaled difference symbols differs from the corresponding
derivative symbol by at most
$C_rM^{-1}|\eta|^{r+1}$ on this region.  These errors are integrable
through $r=8$.  On the complementary region both covariance
multipliers, times a symbol of order $r$, have integral bounded by
\[
 C_{J,L}\int_{M^{1/100}}^\infty t^{r-65}\,\dd t
 \le C_{J,L}M^{-(64-r)/100};
\]
the same estimate bounds the part outside the rescaled Fourier
square for the continuum kernel.  For $r\le8$ this is
$O(M^{-1/2})$.  Fourier inversion proves the stated uniform
scaled-difference estimate.  Its diagonal case and
$G_{L,\upsilon_{\rm en}}(0)=\log L/(2\pi)$ give
$|\delta_j^{\rm en}|\le C_{J,L}(mL^j)^{-1/2}$.
The mean-value bound for $e^{-a'\delta/2}$, uniformly on $I$,
now proves the displayed summable estimate.  The fixed parameter
$\upsilon_{\rm en}$ can leave every compact interval when $J$
varies, which is why this second argument is not a spread-out
estimate.
\end{proof}

We use the block geometry, analytic point norm, and exponential regulator
of Lemmas~\ref{rg:geometry} and~\ref{rg:regulator}; the activity-weight
base is denoted by $A$, to distinguish it from $A_0$.  Activities have a
collar of width $\rho$ rather than one nearest-neighbor step.  Since
$\rho\ll m_0$, enlarging the fixed connection radii contains these
collars without changing the counting or regulator proofs.  Activities
are real, even, periodic under common translations by $2\pi/\alpha$,
and covariant under block translations and square symmetries.  Bulk norm
parameters can be chosen independently of sufficiently large $\rho$.

\begin{proposition}[Microscopic smallness]\label{wh:uv}
For $a'$ in a fixed compact interval containing $k_0$ or lying above it,
the two ultraviolet integrations have the exact block representation
\begin{equation}\label{wh:initial-activity}
 K_0(X)=\ind_{X=\{B\}}\{s e_B^0/2+z_0c_B^\alpha\}
          +\mathcal R_0(X),\qquad
 z_0=2m_0^2e^{-a'D_{\rm uv}/2}.
\end{equation}
For spread-out $J_\rho$, choose $1/2<\gamma'<\gamma$ and put
$\delta_0=e^{-\gamma'a'v^2/2}$.  For exponentially small $s$,
\begin{equation}\label{wh:uv-sharp}
 \|\mathcal R_0\|\le C\rho^C(|s|+\delta_0)^2,
 \qquad
 \|\partial_s\mathcal R_0\|\le C\rho^C(|s|+\delta_0).
\end{equation}
Thus, uniformly on $|s|\le C'z_0$, these bounds are $o(z_0)$ and
$o(1)$, respectively.  Uniformly towards unbounded $a'$, it suffices
to use the weaker bounds $C\rho^C(s^2+\delta_0)$ and
$C\rho^C(|s|+\delta_0)$, and the designated cosine may be absorbed in
the remainder.  For fixed $J$ these weaker bounds hold with $C_J$ and
any fixed $0<\gamma'<\gamma$ when $a'$ is sufficiently large.
\end{proposition}

\begin{proof}
Write the first Gaussian field as $\eta+\zeta$, where
$\Cov(\eta)=\gamma v^2\Id$ and $\Cov(\zeta)=C_{\rm uv}$.
Conditioning each $\eta_x$ on the comb gives the periodic density
\[
 w(t)=\sum_{q\in\Z}e^{-q^2a'\gamma v^2/2}e^{iq\alpha t}.
\]
Expand in occupied site factors $w(\phi_x+\zeta_x)-1$ and in
occupied oriented-edge factors
\[
 e^{sE_\ell(\phi+\zeta+\eta)/2}-1,
 \qquad E_{(x,y)}(\psi)
       =\frac{(\psi_{x+y}-\psi_x)^2}{2|J|v^2}.
\]
Assign an edge to the block of its base vertex and group all occupied
factors by the enlarged block connectivity.  Gaussian range and the
retained collars give exact factorization between separated components.
This is the initial instance of the algebra in
Proposition~\ref{rg:partition-map}.

Here are estimates justifying the expansion and its derivatives.  An
occupied site costs $\delta_0$ in the analytic point norm; its Fourier
derivatives introduce powers of $\alpha|q|$ and factors
$e^{h_0\alpha|q|}$, absorbed by the gap $\gamma-\gamma'$.
An occupied edge costs $C|s|$, using
$|e^{sE/2}-1|\le C|s|(1+E)e^{cE}$ for a fixed arbitrarily small $c$.
For conditional derivatives one differentiates the shifted Gaussian
density while keeping the constrained lattice value fixed.  At the at
most twice as many touched endpoints as occupied edges, the resulting
weights are bounded by
\[
 C\exp\{h_0|\eta_x|/(\gamma v^2)+h_0^2/(2\gamma v^2)\}.
\]
The conditional normalizer and its inverse have the same estimates,
since $w=1+O(\delta_0)$ uniformly in its argument.  These bounds are
uniform also for the first two conditional moments, with polynomial
factors absorbed by the same exponential slack.

The field energy from a selected set of edges is bounded by a constant
times the regulator gradient energy on their collars.  Indeed, resolve
each $y$ into an axis path of length at most $C|y|$ and apply
Cauchy--Schwarz along the path.  Summing its translates and then its
coefficient $1/(|J|v^2)$ uses
$\sum_{y\in J}|y|^2=4|J|v^2$, and gives a bounded constant.
For the $\zeta$ energy, put $c=(|J|v^2)^{-1}$ and let $\mathsf{B}$ have
one row per oriented interaction edge, with
$(\mathsf{B}\psi)_\ell=\sqrt{c/2}\,\nabla_\ell\psi$.  Both edge orientations
are included, so $\mathsf{B}^*\mathsf{B}=A_0$ and
$E_\ell(\psi)=|(\mathsf{B}\psi)_\ell|^2$.
For a selected edge set $S$, let $\mathsf{B}_S$ be the row restriction and
$T_S=\mathsf{B}_SC_{\rm uv}\mathsf{B}_S^*$.
Equation~\eqref{wh:uv-energy} implies
\[
 \|\mathsf{B}C_{\rm uv}^{1/2}\|^2
 =\|C_{\rm uv}^{1/2}A_0C_{\rm uv}^{1/2}\|\le1,
 \qquad 0\le T_S\le\Id,\qquad \tr T_S\le|S|.
\]
Gaussian diagonalization therefore gives, for $0\le\tau<1/2$,
\begin{equation}\label{wh:uv-edge-moment}
 \E\exp\left\{\tau\sum_{\ell\in S}E_\ell(\zeta)\right\}
 =\det(\Id-2\tau T_S)^{-1/2}
 \le\exp\left\{\frac{\tau|S|}{1-2\tau}\right\}.
\end{equation}
Here $-\log(1-x)\le x/(1-x)$ bounds each eigenvalue term.
Polynomial occupation factors are absorbed by
\[
 \prod_{\ell\in S}(1+E_\ell)
 \le C_\varepsilon^{|S|}
                  \exp\left\{\varepsilon\sum_{\ell\in S}E_\ell\right\},
\]
with a fixed reserve below the exponent $1/2$.
Use
$E_\ell(\phi+\zeta+\eta)
 \le3E_\ell(\phi)+3E_\ell(\zeta)+3E_\ell(\eta)$.
The background part is placed directly in the output regulator
by the deterministic path estimate; only the selected edge energies
of $\zeta$ are integrated using Equation~\eqref{wh:uv-edge-moment}.
In particular Lemma~\ref{rg:regulator} is not applied to $C_{\rm uv}$.
For the independent field, each touched vertex is incident to at
most $2|J|$ oriented edges, and hence
\[
 \sum_{\ell\in S}E_\ell(\eta)
 \le\frac2{v^2}\sum_{x\text{ touched}}\eta_x^2.
\]
There are at most $2|S|$ touched vertices.
Shifted-lattice Gaussian sums with the displayed derivative weights
and this sufficiently small quadratic growth are bounded by $C$ per
touched vertex, including their normalizers.  Combining these facts
bounds a term with $r$ occupations by the product of its site and edge
costs, times $C^rW_{m_0}^\kappa$.

On one edge, replacing its conditional expectation by the unconditioned
linear term costs $C\rho^C(s^2+|s|\delta_0)$.  The unconditioned
expectation is exactly $sE_\ell(\phi)/2$ plus a constant.  Extract the
sum of these constants by the constant-decoration identity of
Proposition~\ref{rg:partition-map}; its nonconstant linear part is
$s e_B^0/2$.  A single site's fundamental Fourier terms, after the
$\zeta$ integration, have coefficient $2e^{-a'D_{\rm uv}/2}$.
Summing sites in $B$ gives $z_0c_B^\alpha$.  All its higher harmonics
cost at least $C\delta_0^2$ in the compact-$a'$ regime.

For a prescribed connected support the local choices for $r$
occupations are bounded by $(Cm_0^2|J|)^{Cr}$, by encoding the counts
in each block and the sites and edges for each occupation.  At least
one occupation is required for every occupied block.  Since $m_0$ is
polynomial in $\rho$, this pays any fixed prescribed initial activity
weight.  Summing the absolutely convergent expansion with at least two
occupations, and the single-occupation errors just identified, proves
Equation~\eqref{wh:uv-sharp}.  Differentiating the edge factors proves
the asserted derivative bound by the identical occupation sum with
one power of $|s|$ removed.  The weaker bounds follow by putting all
site occupations, including the fundamental one, in the remainder.
For fixed $J$ all counting factors are fixed, so the same proof gives
its asserted version.  Finally, since $\gamma'>1/2$ and
$D_{\rm uv}=v^2(1+o(1))$, the exponent in $\delta_0^2/z_0$ is
strictly negative for large $\rho$; polynomial factors do not affect
the stated $o(z_0)$ conclusion.
\end{proof}

\subsection{The exact local map}

At scale $m=m_{\rm start}L^j$, use either the first shell family
with $m_{\rm start}=m_0$, or the fixed-$J$ entry family with any
sufficiently large integral $m_{\rm start}$.  Write $\lambda_j$ for
its exact fundamental multiplier, namely $\lambda_j^{\rm en}$ in the
entry case.  Split off the two singleton interactions
$t e_B^0/2+zc_B^\alpha$ and denote the measured remainder by $R$.
Assign their complete Gaussian expectations to the containing coarse
singleton, extracting the constant part of the gradient expectation.
The surviving coefficients are exactly $t$ and $\lambda_jz$.
For small remainder supports localize the neutral part to constants
and $e_B^0/2$, and its fundamental charged part to $c_B^\alpha$.
All these are linear operations on the activity being localized.

We spell out the two modifications of the localizations in
Lemma~\ref{rg:localization}.  If $D^X$ is the Hessian on affine fields
in the standard axes, for each $y\in J$ take the orthonormal frame
$(y,R_{\pi/2}y)/|y|$.  In this frame use the two differences with
steps $y,R_{\pi/2}y$, multiply by $m/|y|$, transform the resulting
vector back to the standard axes, and average the resulting quadratic
forms with weights $|y|^2/(4|J|v^2)$.  Average both orientations as
in Lemma~\ref{rg:localization}.  These weights sum to one.  Axis-path
Taylor bounds with nearest second differences control the affine
replacement error by the same $L^{-3}$ estimate, since $|y|\ll m$.
The prescription commutes with reflections as well as rotations;
a reflection interchanges the frame orientations.  After summing
over $X$ at a fixed block, square symmetry makes its Hessian $d\Id$.
For such a Hessian the assigned expression is exactly
\begin{align*}
 &\frac1{2m^2}\sum_{x\in B}\sum_{y\in J}
  \frac{|y|^2}{4|J|v^2}\,d\frac{m^2}{|y|^2}
    \{(\partial_y\phi_x)^2+(\partial_{R_{\pi/2}y}\phi_x)^2\}
       =\frac d2 e_B^0(\phi).
\end{align*}
Thus the localization preserves the actual long-range energy,
including its allocation to fine and coarse blocks.

For a fundamental projection on a small support $X$, evaluate the
Gaussian expectation on constant fields; divide the coefficient by
$|X|$ and assign it times $c_B^\alpha$ to each $B\in X$.
Evenness cancels the sine coefficients.  For
$\sqrt{k_0}\le\alpha\le2\sqrt{k_0}$, perform the charged contour shift
of Lemma~\ref{rg:charged} with
$i\operatorname{sgn}(q)\alpha\Gamma_m(\cdot,v_0)$.
The density cost and constant character together give
\begin{equation}\label{wh:charge-shift}
 \exp\{-\alpha^2(|q|-1/2)\Gamma_m(0,0)\}.
\end{equation}
Positive-order kernel bounds show that the nonconstant covariance
column has bounded fine test norm independently of $L$.  Choose $h_0$
large enough to contain this residual with strict radius reserve.
The coarse constant derivative factors $e^{h_0\alpha|q|}$ are then
paid by Equation~\eqref{wh:charge-shift} for $L$ large.  Every
$|q|\ge2$ has spare powers of $L$.  For $|q|=1$ the gain is exactly
$\lambda_jL^{-2}$, where $\lambda_j$ is the fundamental multiplier
of the chosen shell family.  It is at most $CL^{-2}$ on a fixed
one-sided interval $k_0\le a'\le k_0+\delta_{\rm ref}$ with
$0<\delta_{\rm ref}<3k_0$ fixed and sufficiently small, once the
summable diagonal errors are small.  Factoring its character at an
anchor in $B$, subtraction
of the constant-background coefficient leaves a first Taylor
difference.  On the small support its nonconstant argument and each
coarse direction have fine test size
$CL^{-1}(1+\sqrt{U_{Lm}})$ and $CL^{-1}$, respectively.  The unused
analytic radius and the small-support regulator slack absorb this
fixed polynomial, giving the further factor $L^{-1}$.  Averaging the
anchor over $B$ proves the claimed localization with $c_B^\alpha$.
By Lemma~\ref{wh:kernels} the diagonal error is arbitrarily small by
taking $\rho$ large, or by increasing the starting scale at fixed $J$.

\begin{proposition}[Weak local maps]\label{wh:map}
Choose the regulator parameters and $h_0$ first, then $L$ sufficiently
large, and then $A$ sufficiently large depending on $L$.  On a fixed
small activity ball the exact partition algebra induces analytic maps
\begin{equation}\label{wh:triangular}
 \begin{pmatrix}t_+\\z_+\\R_+\end{pmatrix}
 =\begin{pmatrix}t+P_j^tR\\\lambda_jz+P_j^zR\\C_jR\end{pmatrix}
     +\mathcal N_j(t,z,R),
 \quad
 \|C_j\|\le\theta<1,\quad
 \|P_j^t\|+\|P_j^z\|\le C_{L,A},
\end{equation}
where $\theta$ can be chosen arbitrarily small and
\begin{equation}\label{wh:nonlinear}
 \|\mathcal N_j(t,z,R)\|
       \le C_{L,A}(|t|+|z|+\|R\|)^2.
\end{equation}
The same ball has uniform analytic Taylor coefficient and differentiated
estimates at fixed $a'$.  These assertions hold on a fixed one-sided
reference interval $k_0\le a'\le k_0+\delta_{\rm ref}$,
uniformly in sufficiently spread-out $J_\rho$ for the first kernel
family.  For every fixed $J$, they hold at a sufficiently large
starting scale for either the first family or the entry family of
Lemma~\ref{wh:kernels}.

On a compact subinterval of $a'>k_0$, one may omit the cosine
coordinate and leave every charged activity in the remainder.  Choose
$L$ depending on the gap from $k_0$; then the corresponding map has
the same assertions with its $z$ row deleted.  For $a'\ge4k_0$ there
is also such a map with a charged contraction uniform in $a'$.

The scalar coefficients obtained on the plane are exactly the torus
coefficients through the stopping scale.  Torus remainders and their
activity derivatives obey the same estimates.  Each fixed-support
finite recursion and each finite scalar projection is continuous in
the parameters while its preceding activities remain in the small
ball.  This last assertion does not assert norm analyticity under
variation of $\alpha$.
\end{proposition}

\begin{proof}
The preceding neutral and fundamental subtractions give
$CL^{-3}$ in point norm before support counting.  There are
$O(L^2)$ small fine supports above a prescribed coarse block, with
weight ratios bounded independently of $L,A$.  Thus these subtracted
linear terms have norm at most $C/L$.  Large supports have the
closure-size gain of Lemma~\ref{rg:geometry}, which makes their
linear contribution arbitrarily small by increasing $A$ after $L$.
The extracted scalar projections are bounded by $C_{L,A}$.
This proves the triangular linear part and $\|C_j\|\le\theta$.

Apply Proposition~\ref{rg:partition-map} with the enlarged collars.
Temporary connectivity contains the covariance range, so convolution
factorizes; all exclusions on an original support are retained when
a localized piece is relocated to a singleton.  Only small supports
shrink, and their coarse diameter is at most one; hence separated
final components carry no remaining mutual exclusion.  Constant
decorations are removed by their convergent inverse series.
The counting and product estimates in that proposition bound every
term with at least two inputs by
$C_{L,A}(|t|+|z|+\|R\|)^2$.  They also converge absolutely for complex
activities in the same real-field norm.  Bounded multilinear sums and
the constant inverse series prove analyticity and the differentiated
bounds, without an appeal to analytic dependence on the period.

If $a'>k_0$ stays in a compact interval, the fundamental charged
estimate before counting is $CL^{-a'/(4\pi)}$.  Its product with the
$O(L^2)$ small-support count is uniformly small once $L$ is chosen
for the positive gap, so fundamental localization is unnecessary.
For $\alpha\ge2\sqrt{k_0}$ use a contour shift of the fixed size
$2\sqrt{k_0}$.  Its gain is
$\exp\{-[2\sqrt{k_0}\alpha|q|-2k_0]\Gamma_m(0,0)\}$.
The coefficient is already at least $2k_0$ for $|q|=1$ and increases
linearly with $\alpha|q|$.  Taking $L$ large pays the output constant
derivative factors uniformly and leaves at least $L^{-4}$ in point
norm.  This proves both no-cosine variants.

A small output support depends on predecessor data within a fixed
number of its blocks: convolution uses a fixed enlargement,
relocations originate only from small supports, and normalization
adds singleton decorations.  Iterating these dependencies contributes
a geometric sum of smaller physical radii, bounded by a fixed
multiple of the current scale.  With $D_0$ larger than this multiple,
the needed covariance entries and activities unwrap identically to
the plane.  Large torus supports retain the closure-size contraction
and need not be identified with plane supports.  Applying the same
norm estimates inductively proves the torus assertion, and applying
their differentiated forms proves its derivative assertion.
At fixed depth, absolute expansions or dominated Gaussian integrals
with regulator reserve give continuity in all parameters of every
fixed local coefficient.  Uniform geometric tails then give
continuity of the scalar projection series used below.
\end{proof}

\subsection{Balancing and the quadratic coefficients}

Let $0\le\mu=1-\lambda\le\mu_0$, with $\mu_0$ small enough that
$\lambda>\theta$.  In this subsection $\lambda_j$ denotes the exact
multiplier of the chosen family: Equation~\eqref{wh:eigenvalue} for
the first family, and $\lambda_j^{\rm en}$ for the entry family.
The operators $C_j,P_j^t,P_j^z$ are likewise those of that family's
map.  Define
\[
 \ell_j=\prod_{h\ge j}\frac{\lambda_h}{\lambda},\qquad
 C_{h:j}=C_{h-1}\cdots C_j\quad(h>j),\qquad C_{j:j}=\Id.
\]
The corresponding summable diagonal estimate makes the product convergent and
$\ell_j=1+o(1)$ uniformly.  Set
\begin{align}
 T&=t+\sum_{h\ge j}P_h^tC_{h:j}R,
       \label{wh:T-balance}\\
 Z&=\ell_jz+
       \sum_{h\ge j}\lambda^{-(h+1-j)}
                         \ell_{h+1}P_h^zC_{h:j}R.
       \label{wh:Z-balance}
\end{align}
Both series converge in operator norm, since their tails are bounded
by $C\theta^r$ and $C(\theta/\lambda)^r$.  The identities
$\ell_{j+1}\lambda_j=\lambda\ell_j$ and cancellation of the first
term in each series show directly that the linear part is
$(T,Z,R)\mapsto(T,\lambda Z,C_jR)$.

\begin{proposition}[Balanced map]\label{wh:balanced-map}
On a fixed small norm ball the maps in these coordinates have the form
\begin{align}
 T_+&=T+b_j^{tt}T^2+b_j^{zz}Z^2+E_{1,j},\notag\\
 Z_+&=(1-\mu)Z+b_j^{tz}TZ+E_{2,j},\label{wh:balanced-equations}\\
 R_+&=C_jR+O((|T|+|Z|+\|R\|)^2),\notag\\
 |E_{i,j}|&\le C\{(|T|+|Z|)^3+(|T|+|Z|)\|R\|+\|R\|^2\}.
 \notag
\end{align}
All errors are analytic in the activity coordinates with the
corresponding Taylor and derivative estimates.  Under half-period
translation the scalar coordinates have parities $T\mapsto T$ and
$Z\mapsto-Z$; this accounts for all displayed quadratic monomials
on pure scalar inputs.
\end{proposition}

\begin{proof}
The linear changes in Equations~\eqref{wh:T-balance} and
\eqref{wh:Z-balance}, together with their inverses, are uniformly
bounded on the product space.  Apply them to the analytic expansion
in Proposition~\ref{wh:map}, keep the quadratic terms containing
only $T,Z$, and place every term with a factor $R$ in the displayed
error.  The quadratic terms in $R_+$ have the same bound by the
inverse coordinate change.  Translation by $\pi/\alpha$ commutes
with convolution, reblocking, and localization, preserves energy,
and changes the sign of the fundamental cosine.  The projection
series commute with this involution as well.  Its even and odd
scalar parities exclude $TZ$ from the first equation and $T^2,Z^2$
from the second.  Cauchy estimates on a smaller fixed ball give all
claimed differentiated bounds.
\end{proof}

\begin{proposition}[The secular quadratic terms]\label{wh:quadratic-coefficients}
In the continuum limit at $a'=k_0$, the balanced coefficients are
independent of the scale and satisfy
\begin{equation}\label{wh:coefficient-signs}
 b^{tt}=0,\qquad b^{tz}=-b<0,\qquad b^{zz}=-H<0.
\end{equation}
The positive numbers $b,H$ stay in compact subintervals of
$(0,\infty)$ as $u$ ranges over the compact interval in
Lemma~\ref{wh:kernels}.  Consequently, given $\varepsilon>0$, one
can first take $\mu_0$ small and then $\rho$ large so that, uniformly
in $j$, $|b_j^{tt}|\le\varepsilon$ and $-b_j^{tz},-b_j^{zz}$ lie in
fixed compact positive intervals.  For every fixed $J$ the same
conclusion follows by making the starting scale sufficiently large.
At that fixed $J$ the coefficients converge to their continuum
limits as $j\to\infty$ when $a'=k_0$.

For the fixed-$J$ entry family, use its own multipliers
$\lambda_j^{\rm en}$ and its own projection operators in the same
balanced coordinates.  Its critical continuum coefficients are
\[
\begin{gathered}
 b_{\rm en}^{tt}=0,\qquad b_{\rm en}^{tz}=-b_{\rm en},
 \qquad b_{\rm en}^{zz}=-H_{\rm en},\\
 b_{\rm en}=2\log L,\qquad
 H_{\rm en}=\frac{8\pi\cdot\pi}{4}
       e^{8\pi c_{\upsilon_{\rm en}}}\log L>0,\\
 \upsilon_{\rm en}=\frac{v}{\sqrt{32}}.
\end{gathered}
\]
At this fixed $J$, the corresponding discrete coefficients converge
to these limits at criticality and have the stated positive bounds
on a sufficiently small one-sided reference interval once $m$ is
large.  If $d$ is the largest power of two not exceeding the sup
range of $J$, then $u_{\rm wh}=v/(d\sqrt{32})$ is the width of
the first fixed-$J$ family at the same block ratio, and
$H_{\rm en}=d^4H_{\rm wh}$, whereas $b_{\rm en}=b_{\rm wh}$.
These entry-family bounds are not uniform in a varying $J$.
\end{proposition}

\begin{proof}
We first justify the limiting calculation, and only then calculate
its coefficients.  A quadratic coefficient in the balanced map is
the quadratic image of its two singleton inputs followed by the
linear projection series.  Truncating those series at depth $r$
costs at most $C\theta^r$ (or $C(\theta/\lambda)^r$).  At any fixed
depth there are finitely many connected block configurations:
quadratic images of pure singleton inputs have bounded size, and
each step has bounded range in the rescaled block coordinates.
Each coefficient is thus a finite volume-normalized sum of Gaussian
polynomial/character expectations and their Taylor projections on
constant or affine backgrounds.  Lemma~\ref{wh:kernels} gives
convergence of every covariance and derivative entering these
expectations.  The rescaled difference in direction $y$, divided
by $|y|$, converges to its directional derivative; the weighted sum
of its squares converges to $\int_B|\nabla\phi|^2$.  The frame
prescription has the same limit even before its Hessian is scalar,
because two frame derivatives transformed back to the axes recover
the entire usual gradient.  Spatial Riemann sums and Gaussian
moment formulas therefore give convergence at finite depth.
Uniform geometric tails permit sending $r\to\infty$ afterwards.
This proves existence of the limiting quadratic map.  Its kernels
are related by dilation, so its coefficients are independent of
scale, and $\ell_j=\lambda=1$ at marginality.

To identify the coefficients, start at side one with the exponential
interaction of $t e_B^0/2+zc_B^\alpha$ and iterate $r$ shells, with
$M=L^r$.  At order one its remainder is zero, and its order-two
initial remainder is bounded independently of $r$.  At order two
the marginal increments add, while
$R^{(2)}_{j+1}=C_jR^{(2)}_j+O(1)$ keeps the remainder coefficients
uniformly bounded.  Thus the final scalar quadratic coefficients
are $r$ times the unknown coefficients in
Equation~\eqref{wh:coefficient-signs}, up to bounded errors.  The
unbalanced coordinates have the same leading terms because their
difference from the balanced coordinates is a bounded projection
of $R^{(2)}$.  This argument is a finite-order calculation in the
exact block algebra; no infinite-volume interacting measure is
assumed to exist for nonzero formal $t,z$.

Two diagnostics read these scalar terms.  The first is the
fundamental cosine coefficient of the connected log density per
output block on a constant background.  The second is its neutral
Hessian in $g_1$ on the affine background
$\theta+g\cdot x/M$, after averaging $\theta$ over one period.
Here connected log density means only its Taylor polynomial through
degree two, computed in a large box and then per output block with
boundary terms removed.  At fixed $r$ all its kernels have bounded
range.  The linear singleton terms read $z$ and $t$ exactly; the
linear remainder terms are bounded by the activity norm.  The
degree-two connected terms of the unchanged first-order singleton
inputs are also bounded in output block units.  An affine
derivative may be anchored in its cluster because of the neutral
phase average, and the resulting fixed powers of its diameter are
summable against its exponential support weight.  Therefore each
diagnostic differs from its output scalar coefficient by $O(1)$,
uniformly in $r$.  Extracted constants disappear from both
diagnostics.  These facts justify reading the secular coefficients
from the original Gaussian integral rather than from its block
representation.

That integral uses the covariance $G_M$ in
Equation~\eqref{wh:continuum-symbol}, with
$G_M(0)=\log M/(2\pi)$.  In the $t^2$ contribution to the neutral
gradient Hessian, a background gradient can occur only in a linear
gradient fluctuation.  Its covariance integrates to a derivative
integral of $G_M$, hence to zero.  Boundary terms vanish on taking
the per-block limit because this kernel has finite range.  There
is no $t^2$ secular term.

For the mixed term, Gaussian differentiation gives
\begin{align*}
 &\Cov\left(\tfrac12|\nabla\zeta(x)|^2,
              \cos(\alpha(\theta+\zeta(y)))\right)\\
 &\hspace{20mm}=-\frac{\alpha^2}{2}|\nabla G_M(x-y)|^2
       e^{-\alpha^2G_M(0)/2}\cos(\alpha\theta).
\end{align*}
The output block area therefore makes the mixed diagnostic
\begin{equation}\label{wh:mixed-diagnostic}
 -\frac{\alpha^2}{2}M^2e^{-\alpha^2G_M(0)/2}
               \int_{\R^2}|\nabla G_M(x)|^2\dd x.
\end{equation}
Parseval and radial integration show
\[
 \int|\nabla G_M|^2
   =\frac1{2\pi}\int_0^\infty
                  (P_u(r)-P_u(Mr))^2\frac{\dd r}{r}
   =\frac{\log M}{2\pi}+O(1).
\]
For example, subtract the indicator of $[M^{-1},1]$ in this last
integral; near each endpoint its error is integrable uniformly by
$P_u(r)=1+O(r^2)$ and $P_u(r)=O(r^{-64})$.
At $\alpha^2=8\pi$, $M^2e^{-\alpha^2G_M(0)/2}=1$.
Dividing Equation~\eqref{wh:mixed-diagnostic} by
$r=\log M/\log L$ gives
\begin{equation}\label{wh:b-value}
 b=\frac{\alpha^2}{4\pi}\log L=2\log L>0.
\end{equation}

For the $z^2$ term, averaging the phase in the connected product of
two cosines gives
\[
 \frac12e^{-\alpha^2G_M(0)}
          (e^{\alpha^2G_M(x-y)}-1)
          \cos\big(\alpha g\cdot(x-y)/M\big).
\]
The second-cumulant factor $1/2$ and two derivatives in $g_1$,
with output area $M^2$ cancelling the affine scaling, give
\begin{equation}\label{wh:neutral-diagnostic}
 -\frac{\alpha^2}{4}e^{-\alpha^2G_M(0)}
       \int_{\R^2}x_1^2(e^{\alpha^2G_M(x)}-1)\dd x.
\end{equation}
We verify its logarithmic asymptotics, including its sign.  Put
\[
 F_u(x)=\int_{\R^2}
       \frac{1-\cos(\eta\cdot x)}{|\eta|^2}
                          P_u(\eta)\frac{\dd\eta}{(2\pi)^2}.
\]
Radial integration gives
$F_u(x)=(2\pi)^{-1}\log|x|+d_u+o(1)$ as $|x|\to\infty$,
uniformly for $u$ in a compact interval.  One direct proof replaces
$P_u$ by the indicator of a unit radial disk: their difference is
integrable in $\dd r/r$ because it is $O(r^2)$ at zero and decays
at infinity.  In the disk term the circle character average at
argument $t$ is $O(t^{-1/2})$.  This bound follows by excluding
arcs of size $t^{-1/2}$ about its two stationary angles and
integrating by parts on the remaining arcs.  Its radial integral
therefore has a convergent constant after subtraction of $\log|x|$.
The same proof yields
$F_u(x)=(2\pi)^{-1}\log(1+|x|)+O(1)$ globally.

Since
$G_M(x)=G_M(0)-F_u(x)+F_u(x/M)$, writing $c_u=-d_u$ gives
\begin{equation}\label{wh:intermediate-log}
 G_M(x)=\frac1{2\pi}\log\frac M{|x|}+c_u+o(1)
       \quad(1\ll|x|\ll M),
 \qquad
 G_M(x)\le\frac1{2\pi}\log^+\frac M{1+|x|}+C.
\end{equation}
For completeness the finite range and logarithmic error also have an
elementary radial representation.  The function
$\sin(u|\eta|)/(u|\eta|)$ is the characteristic function of the
planar projection of the uniform probability measure on the sphere
of radius $u$ in $\R^3$: rotate $\eta$ to an axis and integrate the
uniform axial coordinate on $[-u,u]$.  Thus $P_u$ is the Fourier
transform of its $64$-fold convolution, a radial probability density
$q_u$ supported in the disk of radius $R_u=64u$.  It is bounded and
smooth to every order used here, as follows from the integrability
of $|\eta|^rP_u(\eta)$ for these orders.  The fundamental-solution
identity $\Delta\log|x|=2\pi\delta_0$ gives
\[
 F_u(x)=\frac1{2\pi}\int q_u(y)
                \{\log|x-y|-\log|y|\}\dd y.
\]
Both sides vanish at zero and have Laplacian $q_u$; the radial
harmonic difference is therefore zero.  Angular averaging of the
logarithm is $\log\max\{|x|,|y|\}$, by its mean-value identity
inside or outside the corresponding circle.  Consequently
$F_u(x)=(2\pi)^{-1}\log|x|-c_u$ exactly for $|x|\ge R_u$,
where $c_u=(2\pi)^{-1}\int q_u(y)\log|y|\dd y$.
The scaling $q_u(y)=u^{-2}q_1(y/u)$ also proves
$c_u=c_1+(\log u)/(2\pi)$.  The formula for $G_M$ now shows
$G_M(x)=0$ for $|x|\ge MR_u$, confirming the claimed range without
inferring pointwise positivity from spectral positivity.
At $\alpha^2=8\pi$, the integrand in
Equation~\eqref{wh:neutral-diagnostic} has prefactor $M^{-4}$.
On intermediate logarithmic shells its radial part is
$e^{8\pi c_u}\dd|x|/|x|$, and
$\int_0^{2\pi}\cos^2\vartheta\dd\vartheta=\pi$.
The regions $|x|\le R_u$ and the $-1$ term contribute $O(1)$,
uniformly on compact $u$ intervals.  On $R_u\le|x|\le MR_u$ the
exact formula for $F_u$ makes the exponential part of the radial
integral
\[
 \pi e^{8\pi c_u}\int_{R_u}^{MR_u}
                   e^{8\pi F_u(r/M)}\frac{\dd r}{r}
 =\pi e^{8\pi c_u}\log M+O(1).
\]
Indeed $F_u(t)=O(t^2)$ near zero, and after $t=r/M$ the integral
of $(e^{8\pi F_u(t)}-1)/t$ is bounded on $[0,R_u]$.
It follows that
\begin{equation}\label{wh:H-value}
 H=\frac{\alpha^2\pi}{4}e^{8\pi c_u}\log L>0.
\end{equation}
The continuity and boundedness of $c_u$ on the compact $u$ interval
make this positivity uniform.  The finite-depth convergence proved
at the start, first uniformly in small $\mu$ and then in the
summable lattice error, now proves all discrete assertions of the
proposition.  This computes both signs directly in the compensated
positive-gradient convention of Equation~\eqref{wh:compensation}.

For the entry family, the fixed-depth Gaussian expectations at the
start of the proof converge by its order-eight scaled-difference
bound, and the same geometric projection tails are uniform once its
own diagonal errors are small.  The localizations and normalized
singleton inputs are unchanged: in particular
$c_B^\alpha=M^{-2}\sum_{x\in B}\cos(\alpha\phi_x)$ converges to the
unit-block cosine integral in both families.  The limiting
calculation therefore uses $G_{M,\upsilon_{\rm en}}$ in the two
diagnostics.  Equations~\eqref{wh:b-value} and~\eqref{wh:H-value}
give precisely $b_{\rm en}$ and $H_{\rm en}$ displayed above.
The identity $c_u=c_1+(\log u)/(2\pi)$ gives
$e^{8\pi(c_{\upsilon_{\rm en}}-c_{u_{\rm wh}})}=d^4$.
No coordinate factor changes this number: the balanced factor
$\ell_j$ tends to one at criticality in either continuum limit.
\end{proof}

\subsection{Shooting away from marginality}

\begin{proposition}[A vanishing-activity shot]\label{wh:gapped-shooting}
In each of the following regimes there is a choice of the input
$s$ in Equation~\eqref{wh:compensation} for which all interaction
activities tend to zero exponentially in scale:
\begin{enumerate}[label=(\roman*)]
 \item $J=J_\rho$, a prescribed fixed $a'>k_0$, and $\rho$
       sufficiently large;
 \item an arbitrary fixed admissible $J$ and $a'$ sufficiently large.
\end{enumerate}
The shot satisfies $|s|\le Cx_*$, where $x_*$ is exponentially
small in $a'v^2$ in (i), and $x_*\le C_Je^{-c_Ja'}$ in (ii).
At the input $k=a'/(1-s)$ its actual structural coefficient is
exactly $a'$.
\end{proposition}

\begin{proof}
Use the no-cosine map of Proposition~\ref{wh:map}.  In (i) take
$L$ for the gap $a'-k_0$ and then $\rho$ large; in (ii) use its
uniform high-frequency contraction.  Balance $t$ by
Equation~\eqref{wh:T-balance}, now with no $Z$ coordinate.  The
recursion is
\[
 T_{j+1}=T_j+O((|T_j|+\|R_j\|)^2),\qquad
 R_{j+1}=C_jR_j+O((|T_j|+\|R_j\|)^2).
\]
Proposition~\ref{wh:uv}, after allowing exponent slack, supplies
$\|R_0\|\le Cx_*$ and $T_0=s+O(x_*)$ on
$|s|\le C'x_*$.  The constant in the last $O(x_*)$ can be made
independent of any fixed $C'$: the dependence on $s$ outside its
linear term is at most $C(s^2+|s|\delta_0)$ and hence is
$o(x_*)$ on every such interval.

Choose $\theta<\theta_3<1$ and a fixed large envelope constant
$D$.  While
\begin{equation}\label{wh:shoot-envelope}
 |T_j|\le Dx_*\theta_3^j
\end{equation}
holds, the remainder recursion and geometric summation imply
$\|R_j\|\le D_Rx_*\theta_3^j$ for a fixed $D_R$:
its forced contribution is bounded by
$Cx_*^2\sum_{h<j}\theta^{j-1-h}\theta_3^{2h}$, which is at
most $C'x_*^2\theta_3^j$.  This estimate also propagates torus
remainders through their stopping scales.  Choose $C'$ so that
the two endpoints $s=\pm C'x_*$ start outside the two opposite
sides of Equation~\eqref{wh:shoot-envelope}.

If equality first occurs at scale $j$, then
$|T_{j+1}-T_j|\le Cx_*^2\theta_3^{2j}$, whereas the envelope
decreases by $D(1-\theta_3)x_*\theta_3^j$.  For small enough
$x_*$ the next iterate exits strictly on the same side.  A strict
exit persists under a sufficiently small change of $s$ by
finite-time continuity.  Thus the sets of parameters eventually
exiting on the positive and on the negative side are disjoint
relatively open sets; when equality occurs its strict next exit
gives the same openness.  Both sets are nonempty.  If there were
no survivor, they would disconnect the shooting interval, which
is impossible.  The maps can be evaluated for the one extra step
at an exit because the whole enlarged envelope lies in their
fixed small domain.  A survivor obeys
Equation~\eqref{wh:shoot-envelope} and the remainder estimate at
all scales, proving decay of every activity.

It remains to identify the shot with the physical coefficient.
For a smooth mean-zero continuum test $f$, let $f_n$ be its sampled
lattice source including the area $n^{-2}$ and discrete mean
subtraction.  Completing the Gaussian square in the original
comb integral gives the prefactor
$\exp\{(f_n,A_0^{-1}f_n)/2\}$ for the $\phi$ field and the ratio
of the full interaction integrals at the translated background
$F_n=A_0^{-1}f_n$ and at zero.  This translation includes the
constraint itself before ultraviolet integration.  Therefore all
field-independent constants extracted at every scale cancel
exactly in the ratio.  At $m_*$ there are boundedly many blocks.
The terminal field obeys the regulator estimates with reserve;
$F_n$ and its first two scaled differences are uniformly bounded
by discrete Fourier summation of the smooth sampled source.
The activity decay thus makes each of the two terminal interaction
averages tend to one, by the finite block expansion and dominated
Gaussian integration.  Finally,
\[
 (f_n,A_0^{-1}f_n)\longrightarrow
       (f,(-\Delta_{\T})^{-1}f)_{L^2(\T^2)}
\]
by the quadratic symbol at zero and the rapid Fourier decay of
$f$.  Since $\sigma=\sqrt{a'}\phi$, the height Laplace limit has
coefficient $a'$.  We next prove positivity of the structural
coefficient at this input using only its unconditional free-rectangle
upper bound, before using the positive-coefficient transfer theorem.

Fix a nonzero smooth real mean-zero $f$ and let
$g=(-\Delta_{\T^2})^{-1}f$, with zero mean.  Write $g_n$ for its
sample and $\Var_{\T_n}$ for variance in the physical height law
on the torus of side $n$.  For a fixed integer $q\ge2$, partition a
fundamental square into $q^2$ axis rectangles by rounding the cuts
$in/q$, $0\le i\le q$, and delete every interaction edge joining
different rectangles or crossing the fundamental-square seam.
For all sufficiently large $n$, these are free rectangles of
side comparable to $n/q$.  Define
\[
 X_n=(\sigma,A_0g_n),\qquad
 X_n^{\rm cut}=\sum_{B\text{ in the partition}}
                      (\sigma|_B,A_Bg_n|_B).
\]
There are at most $C_{J,q}n$ deleted edges, and periodic smoothness
gives $|\nabla_e g_n|\le C_{J,g}/n$ on all of them, including
the wraparound edges.  The difference $X_n-X_n^{\rm cut}$ is
represented by these edge flows with weights $c_e\nabla_e g_n$,
so its flow cost in Equation~\eqref{st:flow} is at most
$C_{J,q,g}/n$.  Lemma~\ref{st:finite-gaussian} therefore gives
\begin{equation}\label{wh:seam-error}
 \Var_{\T_n}(X_n-X_n^{\rm cut})\le C_{J,q,g}k/n.
\end{equation}
Each summand in $X_n^{\rm cut}$ annihilates the constant on its
rectangle.  Deleting the crossing precisions and releasing the
rectangle constants are thus permissible in the free-graph
comparison of Lemma~\ref{st:finite-gaussian}.  Under the cut law
the rectangle representatives are independent, and hence
\[
 \Var_{\T_n}(X_n^{\rm cut})
   \le\sum_B\Var_{B,k}(\sigma,A_Bg_n|_B).
\]
For this fixed finite partition, apply the unconditional
Equation~\eqref{st:upper-energy} along $n=L^N\to\infty$ to
each sequence of rectangles.  Their rounded endpoints converge
to the indicated continuum endpoints.  The $L^2$ triangle
inequality and Equation~\eqref{wh:seam-error} give
\[
 \sqrt{\Var_{\T_n}X_n}
 \le\left\{\sum_B\Var_{B,k}(\sigma,A_Bg_n|_B)\right\}^{1/2}
       +C_{J,q,g}\sqrt{k/n},
\]
and therefore
\begin{equation}\label{wh:unconditional-torus-upper}
 \limsup_{N\to\infty}\Var_{\T_{L^N}}(\sigma,A_0g_{L^N})
      \le a(k)\sum_B\int_B|\nabla g|^2
      =a(k)(f,(-\Delta_{\T^2})^{-1}f).
\end{equation}
No positivity assumption on $a(k)$ enters this comparison, and
no limit in $q$ is required.

To replace the gradient observation by the density source used
in the shot, put $e_n=f_n-A_0g_n$.  The quadratic symbol and
smooth finite differences imply
$n^{-2}\sum_x|n^2e_n(x)|^2\to0$, with $\sum_xe_n(x)=0$.
The torus Poincar\'e bound from the nearest-neighbor edges gives
\[
 \Var_{\T_n}(\sigma,e_n)
 \le k(e_n,A_0^+e_n)
 \le C_Jk n^2\sum_x|e_n(x)|^2\longrightarrow0.
\]
Thus Equation~\eqref{wh:unconditional-torus-upper} holds with
$f_n$ in place of $A_0g_n$.  Bare Gaussian domination for every
fixed real multiple of $f_n$ gives uniform exponential moments;
the shot's Laplace limit consequently gives convergence of its
second moments.  Its limiting variance is
$a'(f,(-\Delta_{\T^2})^{-1}f)$.  The energy of the chosen nonzero
$f$ is positive, so Equation~\eqref{wh:unconditional-torus-upper}
implies $a(k)\ge a'>0$.  Theorem~\ref{an:dirichlet-torus} now
applies on the positive branch and identifies the same torus
Laplace limit with coefficient $a(k)$.  Uniqueness on this one
nonzero test proves $a(k)=a'$, as claimed.
\end{proof}

\begin{corollary}[Every sufficiently high input]\label{wh:high-temperature}
For every admissible fixed $J$, the positive-coefficient regime is
nonempty and its physical onset is finite.  Moreover,
\begin{equation}\label{wh:normalization}
 a(k)=k+O_J(e^{-c_Jk}),\qquad
 \beff(J,\beta)=\beta+O_J(e^{-c_J'\beta})
       \quad(k,\beta\to\infty).
\end{equation}
\end{corollary}

\begin{proof}
For each sufficiently large $a'$, Proposition~\ref{wh:gapped-shooting}
gives an input $k(a')$ with $a(k(a'))=a'$ and
$|k(a')-a'|\le C_Je^{-c_Ja'}$, after absorbing the factor $a'$
into a smaller exponent.  This already proves nonemptiness and
finite onset by the equivalence of positive coefficient and physical
roughness in Corollary~\ref{an:physical-criterion}.  To cover every high input $k$, choose
$\delta_k=e^{-c_Jk/2}$ and shoot at $a'_- =k-\delta_k$ and
$a'_+=k+\delta_k$.  For large $k$ the shot errors are less than
$\delta_k/2$, so $k(a'_- )<k<k(a'_+)$.  The monotonicity in
Proposition~\ref{an:coefficient-regularity}
gives $k-\delta_k\le a(k)\le k+\delta_k$.
This proves the first assertion of Equation~\eqref{wh:normalization}.
Substitute $k=\beta/v^2$ and $\beff=v^2a(k)$ to obtain the second.
\end{proof}

\begin{corollary}[Spread-out onset]\label{wh:spread-out}
For the spread-out family,
\[
 \frac{\bc(J_\rho)}{8\pi v_{J_\rho}^2}\longrightarrow1.
\]
\end{corollary}

\begin{proof}
Corollary~\ref{st:charge-gap} and Corollary~\ref{an:physical-criterion} give
$k_c(J_\rho)=\bc(J_\rho)/v_{J_\rho}^2\ge8\pi$.
For every fixed $a'>8\pi$, regime (i) of
Proposition~\ref{wh:gapped-shooting} gives a rough input
$k(a',\rho)=a'+o(1)$.  Hence
$\limsup_{\rho\to\infty}k_c(J_\rho)\le a'$.
Letting this fixed $a'$ decrease to $8\pi$ proves the result.
The block factor may depend on $a'-8\pi$ in this argument; the
limit is taken in $\rho$ before decreasing $a'$, so that dependence
does not affect the squeeze.
\end{proof}
